% =============================================================================
%  progress-report-01.tex — รายงานความก้าวหน้าโครงการ ULMs ครั้งที่ 1
%  *** ฉบับนี้คือฉบับที่ส่งอาจารย์จริง ***
%
%  build:  cd ~/Projects/SoftwareEn/docs && latexmk progress-report-01.tex
%          ผลลัพธ์ออกที่ build/progress-report-01.pdf
%
%  หลักการเขียนฉบับนี้ — อ่านก่อนแก้
%    1) สั้น อาจารย์มีเวลาจำกัด เป้าหมายคือไม่เกิน 12 หน้า ไม่มีสารบัญ
%    2) มี "แผนปัจจุบัน" เส้นเดียวเท่านั้น ห้ามพล็อตแผนก่อนปรับปรุงลงในกราฟ
%    3) ห้ามอ้างถึงรายงานฉบับก่อน ๆ ทุกกรณี — สำหรับอาจารย์นี่คือรายงานฉบับแรก
%       (ฉบับร่างภายใน progress.tex / progress2.tex / progress3.tex ไม่เคยส่ง)
%    4) รายงานความคืบหน้าที่เกิดขึ้นแล้ว ไม่มีข้อเสนอขอตัดขอบเขต
%    5) ตัดหัวข้อที่ไม่จำเป็นต่อการประเมินความคืบหน้าออก ได้แก่ ปัญหาอุปสรรค
%       การประเมินความเสี่ยง ทรัพยากรที่ต้องการเพิ่ม และภาคผนวกที่มาของตัวเลข
%       ข้อจำกัดที่ตรวจพบยังคงอยู่ แต่ย้ายไปอยู่ใต้ผลงานแต่ละฝั่ง
%    6) ห้ามรายงานผลงานเป็นจำนวนบรรทัดโค้ดหรือจำนวนหน้าเอกสาร ตัววัดพวกนี้
%       ไม่สะท้อนความก้าวหน้าจริง ให้บอกว่า "ทำอะไรได้แล้ว" แทน
%    7) อุปกรณ์ระดับ T3 (ของติดที่ / การจองสถานที่) อยู่ในขอบเขตที่ต้องทำ
%       ไม่ใช่งานที่ตัดออก ต้องปรากฏในแผน Sprint
%    8) เนื้อความ (ย่อหน้าและรายการ) ใช้ตัวอักษรน้ำหนักเดียวกันทั้งหมด
%       ห้ามใช้ \textbf เน้นคำกลางย่อหน้า ตัวหนาสงวนไว้ให้หัวข้อ หัวตาราง
%       และหน้าปกเท่านั้น
%
%  ตัวเลขทุกตัวอ้างสถานะของ origin/main ที่ commit 9b4ab92 ณ 27 ส.ค. 2569
%  ไม่นับโฟลเดอร์ spike-option-b/
%
%  รอบ 27 ส.ค. — นับเฉพาะสิ่งที่รวมเข้า origin/main แล้วเท่านั้น ตามมติของทีม
%  งานบน branch ที่ยังไม่ merge (feat/trpc-staff-staff = วงจรยืม-คืนฝั่ง
%  เจ้าหน้าที่, feat/trpc-auth-admin = รายการอุปกรณ์+เครดิต, front-end-branch
%  = หน้ารับของ/ใช้ห้อง) เขียนเสร็จและ CI เขียวแล้ว แต่ไม่ถูกนับในเปอร์เซ็นต์
%  จึงเล่าไว้ในเนื้อความแทน ถ้ารอบหน้าเปลี่ยนวิธีนับ ต้องแก้หมายเหตุนี้ด้วย
%
%  วันรายงานรอบนี้ตรงกับขอบ Sprint 2 พอดี ค่าแผนจึงอ่านจากตารางภาคผนวก
%  แถว S2 ได้ตรง ๆ ที่ 65 ไม่ต้องสอดแทรกเชิงเส้นเหมือนรอบก่อน
% =============================================================================
% =============================================================================
%  ulms-preamble.tex — preamble ที่ใช้ร่วมกันทุกเอกสารของโครงการ ULMs
%
%  ใช้โดย:  main.tex (ข้อเสนอโครงการ) · progress.tex (รายงานความก้าวหน้า)
%
%  เอกสารที่เรียกใช้ต้องทำสองอย่างหลัง % =============================================================================
%  ulms-preamble.tex — preamble ที่ใช้ร่วมกันทุกเอกสารของโครงการ ULMs
%
%  ใช้โดย:  main.tex (ข้อเสนอโครงการ) · progress.tex (รายงานความก้าวหน้า)
%
%  เอกสารที่เรียกใช้ต้องทำสองอย่างหลัง % =============================================================================
%  ulms-preamble.tex — preamble ที่ใช้ร่วมกันทุกเอกสารของโครงการ ULMs
%
%  ใช้โดย:  main.tex (ข้อเสนอโครงการ) · progress.tex (รายงานความก้าวหน้า)
%
%  เอกสารที่เรียกใช้ต้องทำสองอย่างหลัง \input{ulms-preamble}:
%    1) \hypersetup{pdftitle={...}}  — ชื่อเอกสารใน metadata ของ PDF
%    2) \begin{document}
%  (pdftitle ถูกย้ายออกจากไฟล์นี้ เพราะเป็นค่าเฉพาะของแต่ละเอกสาร)
%
%  แก้ที่นี่ที่เดียว = มีผลกับทุกเอกสาร ห้ามก๊อป preamble ไปวางซ้ำ
% =============================================================================
% =============================================================================
%  main.tex  —  เอกสารข้อเสนอโครงการ (Academic Project Proposal)
%  ระบบบริหารจัดการการยืม–คืนอุปกรณ์มหาวิทยาลัย (ULMs)
%
%  IMPORTANT: This document contains Thai text and MUST be compiled with
%  XeLaTeX (not pdflatex). Thai needs a Unicode font + ICU line-breaking.
%      Build:  latexmk main.tex        (uses .latexmkrc -> xelatex)
%      or:     xelatex main.tex        (run twice for the table of contents)
% =============================================================================

\documentclass[11pt, a4paper]{article}

% ---- Fonts & Thai support ----------------------------------------------------
% fontspec lets XeLaTeX use any system font. Laksaman is a complete serif face
% that covers BOTH Thai and Latin glyphs, so one \setmainfont handles the whole
% (Thai + English) document without font switching.
\usepackage{fontspec}

% Thai has no spaces between words, so TeX cannot find line-break points on its
% own. These two XeTeX primitives switch on ICU's dictionary-based Thai word
% segmentation, which is what makes Thai paragraphs wrap correctly.
\XeTeXlinebreaklocale "th"
\XeTeXlinebreakskip = 0pt plus 0.1pt

\setmainfont{Laksaman}

% ---- Page geometry & paragraph style ----------------------------------------
\usepackage[margin=2.5cm]{geometry}
\usepackage{parskip}          % block paragraphs (no indent, gap between) — common
                              % and more readable for Thai documents
\linespread{1.15}             % a little extra leading; Thai has tall tone marks

% ---- Colours (monochrome / grayscale palette) ------------------------------
\usepackage[table]{xcolor}    % 'table' option enables \rowcolor for table headers
% Monochrome (grayscale) palette. Names kept as-is so every reference still
% works; only the tones changed from the original green/blue accents.
\definecolor{kugreen}{HTML}{1A1A1A}   % near-black — section headings & title
\definecolor{kudark}{HTML}{333333}    % dark gray — subsection/subsubsection titles
\definecolor{headbg}{HTML}{EAEAEA}    % light gray background for table header rows
\definecolor{linkblue}{HTML}{404040}  % medium-dark gray — links / citations

% ---- Section heading styling -------------------------------------------------
% 'nobottomtitles*' stops a heading from being printed near the bottom of a
% page: if less than \bottomtitlespace of room is left, the heading (and its
% following text) moves to the next page. This is titlesec's built-in, robust
% replacement for the earlier \needspace hack — the latter injected vertical
% glue that clashed with longtable page-splitting ("Infinite glue shrinkage").
\usepackage[nobottomtitles*]{titlesec}
\titleformat{\section}
  {\color{kugreen}\Large\bfseries}{\thesection}{0.6em}{}
\titleformat{\subsection}
  {\color{kudark}\large\bfseries}{\thesubsection}{0.5em}{}
\titleformat{\subsubsection}
  {\color{kudark}\normalsize\bfseries}{\thesubsubsection}{0.5em}{}
\titlespacing*{\section}{0pt}{1.4\baselineskip}{0.5\baselineskip}

% Minimum room a heading needs at the bottom of a page; if only ~1–3 lines are
% left (less than this), the heading is bumped to the next page instead of being
% stranded above its table/text.
\renewcommand{\bottomtitlespace}{4\baselineskip}

% ---- Hierarchical indentation by heading level ------------------------------
% Each heading level shifts BOTH the heading and the body text under it further
% to the right: section = 0, subsection (x.y) = 1 step, subsubsection (x.y.z) =
% 2 steps. We do this by having the sectioning commands set \leftskip (a global
% left margin that carries through every following line until the next heading
% resets it). \par first, so the *previous* paragraph is not re-indented.
% \leftskip indents plain paragraphs; \@totalleftmargin makes list environments
% (itemize/enumerate) start from the same indent instead of the page margin.
\makeatletter
\newlength{\lvlindent}
\setlength{\lvlindent}{1.8em}          % one indent step
\let\ulmsSecIndent\section
\renewcommand{\section}{\par\global\leftskip=0pt\global\@totalleftmargin=0pt\relax\ulmsSecIndent}
\let\ulmsSubsecIndent\subsection
\renewcommand{\subsection}{\par\global\leftskip=\lvlindent\global\@totalleftmargin=\lvlindent\relax\ulmsSubsecIndent}
\let\ulmsSubsubsecIndent\subsubsection
\renewcommand{\subsubsection}{\par\global\leftskip=2\lvlindent\global\@totalleftmargin=2\lvlindent\relax\ulmsSubsubsecIndent}
\makeatother

% Indent the headings themselves to match (titlesec ignores \leftskip for the
% heading line, so its left margin is set here via the first \titlespacing arg).
\titlespacing*{\subsection}   {\lvlindent} {1.0\baselineskip}{0.4\baselineskip}
\titlespacing*{\subsubsection}{2\lvlindent}{0.8\baselineskip}{0.3\baselineskip}

% ---- Lists (tighter, custom bullets) ----------------------------------------
\usepackage{enumitem}
\setlist[itemize]{leftmargin=1.6em, itemsep=2pt, topsep=3pt}
\setlist[enumerate]{leftmargin=1.9em, itemsep=2pt, topsep=3pt}

% ---- Tables ------------------------------------------------------------------
\usepackage{array}
\usepackage{booktabs}         % \toprule / \midrule / \bottomrule
\usepackage{tabularx}         % auto-width columns that wrap long Thai text
\usepackage{longtable}        % tables that can break across pages
\usepackage{multirow}         % merged cells for the team-structure table

% Left-aligned wrapping column types:
%   L{width} — fixed-width ragged-right paragraph column
%   Y        — flexible ragged-right X column (for tabularx)
\newcolumntype{L}[1]{>{\raggedright\arraybackslash}p{#1}}
\newcolumntype{Y}{>{\raggedright\arraybackslash}X}
\renewcommand{\arraystretch}{1.3}

% Helpers for a styled header row inside tables.
%   \hrow  — starts a header row and shades it (must be FIRST in the row;
%            colortbl allows only one \rowcolor per row).
%   \thead — bold header-cell text.
\newcommand{\hrow}{\rowcolor{headbg}}
\newcommand{\thead}[1]{\textbf{#1}}

% \begin{keeptable} … \end{keeptable} — keeps a table and the bold label above
% it on the SAME page (used by the data dictionary in §5). A tabularx cannot
% split across pages, so LaTeX is free to break between the label and its
% table, stranding the label alone at the bottom of a page; wrapping the pair
% in a minipage makes them one unbreakable block.
%   * \leftskip is zeroed inside: the minipage box is already placed at the
%     current heading indent, so leaving it set would indent the content twice
%     and push the table past the right margin.
%   * \parskip is restored because minipage's \@parboxrestore zeroes it, which
%     would glue the table straight onto the label line.
% Tables inside must use \linewidth (= the minipage width), not \textwidth.
\newenvironment{keeptable}
  {\par\medskip\noindent
   \begin{minipage}{\dimexpr\textwidth-\leftskip\relax}%
   \leftskip=0pt\relax
   \setlength{\parskip}{0.5\baselineskip plus 2pt}}
  {\end{minipage}\par\medskip}

% ---- Table of contents styling ----------------------------------------------
% By default LaTeX gives dotted leaders to subsections but NOT to sections
% (they print bold with just a right-aligned page number). tocloft lets us turn
% on dots for the section level too, so every TOC line runs "….." to its page.
\usepackage{tocloft}
\renewcommand{\cftsecdotsep}{\cftdotsep}   % dotted leaders for section entries
\renewcommand{\cftsecleader}{\cftdotfill{\cftsecdotsep}}
\renewcommand{\cftsecpagefont}{\normalfont}  % keep page numbers un-bold

% ---- Figures -----------------------------------------------------------------
\usepackage{graphicx}
\graphicspath{{figures/}}

% Landscape (rotated) full-page floats — needed by the ER diagram in §5, which
% is far wider than it is tall and would be unreadable scaled to \textwidth.
\usepackage{rotating}

% ---- Flowcharts / diagrams (TikZ) -------------------------------------------
% Used for the borrow–return workflow diagram (§5). Libraries: geometric shapes
% for decision diamonds, arrows.meta for arrowheads, positioning for relative
% node placement, calc/fit/backgrounds for routing loops and grouping.
\usepackage{tikz}
\usetikzlibrary{shapes.geometric, arrows.meta, positioning, calc, fit, backgrounds}

% ---- Links & bookmarks -------------------------------------------------------
\usepackage{hyperref}
\hypersetup{
  colorlinks=true, linkcolor=black, citecolor=linkblue, urlcolor=linkblue,
  pdflang={th},
}

% Thai labels for auto-generated headings.
\renewcommand{\contentsname}{สารบัญ (Table of Contents)}
\renewcommand{\tablename}{ตาราง}
\renewcommand{\figurename}{รูปที่}

:
%    1) \hypersetup{pdftitle={...}}  — ชื่อเอกสารใน metadata ของ PDF
%    2) \begin{document}
%  (pdftitle ถูกย้ายออกจากไฟล์นี้ เพราะเป็นค่าเฉพาะของแต่ละเอกสาร)
%
%  แก้ที่นี่ที่เดียว = มีผลกับทุกเอกสาร ห้ามก๊อป preamble ไปวางซ้ำ
% =============================================================================
% =============================================================================
%  main.tex  —  เอกสารข้อเสนอโครงการ (Academic Project Proposal)
%  ระบบบริหารจัดการการยืม–คืนอุปกรณ์มหาวิทยาลัย (ULMs)
%
%  IMPORTANT: This document contains Thai text and MUST be compiled with
%  XeLaTeX (not pdflatex). Thai needs a Unicode font + ICU line-breaking.
%      Build:  latexmk main.tex        (uses .latexmkrc -> xelatex)
%      or:     xelatex main.tex        (run twice for the table of contents)
% =============================================================================

\documentclass[11pt, a4paper]{article}

% ---- Fonts & Thai support ----------------------------------------------------
% fontspec lets XeLaTeX use any system font. Laksaman is a complete serif face
% that covers BOTH Thai and Latin glyphs, so one \setmainfont handles the whole
% (Thai + English) document without font switching.
\usepackage{fontspec}

% Thai has no spaces between words, so TeX cannot find line-break points on its
% own. These two XeTeX primitives switch on ICU's dictionary-based Thai word
% segmentation, which is what makes Thai paragraphs wrap correctly.
\XeTeXlinebreaklocale "th"
\XeTeXlinebreakskip = 0pt plus 0.1pt

\setmainfont{Laksaman}

% ---- Page geometry & paragraph style ----------------------------------------
\usepackage[margin=2.5cm]{geometry}
\usepackage{parskip}          % block paragraphs (no indent, gap between) — common
                              % and more readable for Thai documents
\linespread{1.15}             % a little extra leading; Thai has tall tone marks

% ---- Colours (monochrome / grayscale palette) ------------------------------
\usepackage[table]{xcolor}    % 'table' option enables \rowcolor for table headers
% Monochrome (grayscale) palette. Names kept as-is so every reference still
% works; only the tones changed from the original green/blue accents.
\definecolor{kugreen}{HTML}{1A1A1A}   % near-black — section headings & title
\definecolor{kudark}{HTML}{333333}    % dark gray — subsection/subsubsection titles
\definecolor{headbg}{HTML}{EAEAEA}    % light gray background for table header rows
\definecolor{linkblue}{HTML}{404040}  % medium-dark gray — links / citations

% ---- Section heading styling -------------------------------------------------
% 'nobottomtitles*' stops a heading from being printed near the bottom of a
% page: if less than \bottomtitlespace of room is left, the heading (and its
% following text) moves to the next page. This is titlesec's built-in, robust
% replacement for the earlier \needspace hack — the latter injected vertical
% glue that clashed with longtable page-splitting ("Infinite glue shrinkage").
\usepackage[nobottomtitles*]{titlesec}
\titleformat{\section}
  {\color{kugreen}\Large\bfseries}{\thesection}{0.6em}{}
\titleformat{\subsection}
  {\color{kudark}\large\bfseries}{\thesubsection}{0.5em}{}
\titleformat{\subsubsection}
  {\color{kudark}\normalsize\bfseries}{\thesubsubsection}{0.5em}{}
\titlespacing*{\section}{0pt}{1.4\baselineskip}{0.5\baselineskip}

% Minimum room a heading needs at the bottom of a page; if only ~1–3 lines are
% left (less than this), the heading is bumped to the next page instead of being
% stranded above its table/text.
\renewcommand{\bottomtitlespace}{4\baselineskip}

% ---- Hierarchical indentation by heading level ------------------------------
% Each heading level shifts BOTH the heading and the body text under it further
% to the right: section = 0, subsection (x.y) = 1 step, subsubsection (x.y.z) =
% 2 steps. We do this by having the sectioning commands set \leftskip (a global
% left margin that carries through every following line until the next heading
% resets it). \par first, so the *previous* paragraph is not re-indented.
% \leftskip indents plain paragraphs; \@totalleftmargin makes list environments
% (itemize/enumerate) start from the same indent instead of the page margin.
\makeatletter
\newlength{\lvlindent}
\setlength{\lvlindent}{1.8em}          % one indent step
\let\ulmsSecIndent\section
\renewcommand{\section}{\par\global\leftskip=0pt\global\@totalleftmargin=0pt\relax\ulmsSecIndent}
\let\ulmsSubsecIndent\subsection
\renewcommand{\subsection}{\par\global\leftskip=\lvlindent\global\@totalleftmargin=\lvlindent\relax\ulmsSubsecIndent}
\let\ulmsSubsubsecIndent\subsubsection
\renewcommand{\subsubsection}{\par\global\leftskip=2\lvlindent\global\@totalleftmargin=2\lvlindent\relax\ulmsSubsubsecIndent}
\makeatother

% Indent the headings themselves to match (titlesec ignores \leftskip for the
% heading line, so its left margin is set here via the first \titlespacing arg).
\titlespacing*{\subsection}   {\lvlindent} {1.0\baselineskip}{0.4\baselineskip}
\titlespacing*{\subsubsection}{2\lvlindent}{0.8\baselineskip}{0.3\baselineskip}

% ---- Lists (tighter, custom bullets) ----------------------------------------
\usepackage{enumitem}
\setlist[itemize]{leftmargin=1.6em, itemsep=2pt, topsep=3pt}
\setlist[enumerate]{leftmargin=1.9em, itemsep=2pt, topsep=3pt}

% ---- Tables ------------------------------------------------------------------
\usepackage{array}
\usepackage{booktabs}         % \toprule / \midrule / \bottomrule
\usepackage{tabularx}         % auto-width columns that wrap long Thai text
\usepackage{longtable}        % tables that can break across pages
\usepackage{multirow}         % merged cells for the team-structure table

% Left-aligned wrapping column types:
%   L{width} — fixed-width ragged-right paragraph column
%   Y        — flexible ragged-right X column (for tabularx)
\newcolumntype{L}[1]{>{\raggedright\arraybackslash}p{#1}}
\newcolumntype{Y}{>{\raggedright\arraybackslash}X}
\renewcommand{\arraystretch}{1.3}

% Helpers for a styled header row inside tables.
%   \hrow  — starts a header row and shades it (must be FIRST in the row;
%            colortbl allows only one \rowcolor per row).
%   \thead — bold header-cell text.
\newcommand{\hrow}{\rowcolor{headbg}}
\newcommand{\thead}[1]{\textbf{#1}}

% \begin{keeptable} … \end{keeptable} — keeps a table and the bold label above
% it on the SAME page (used by the data dictionary in §5). A tabularx cannot
% split across pages, so LaTeX is free to break between the label and its
% table, stranding the label alone at the bottom of a page; wrapping the pair
% in a minipage makes them one unbreakable block.
%   * \leftskip is zeroed inside: the minipage box is already placed at the
%     current heading indent, so leaving it set would indent the content twice
%     and push the table past the right margin.
%   * \parskip is restored because minipage's \@parboxrestore zeroes it, which
%     would glue the table straight onto the label line.
% Tables inside must use \linewidth (= the minipage width), not \textwidth.
\newenvironment{keeptable}
  {\par\medskip\noindent
   \begin{minipage}{\dimexpr\textwidth-\leftskip\relax}%
   \leftskip=0pt\relax
   \setlength{\parskip}{0.5\baselineskip plus 2pt}}
  {\end{minipage}\par\medskip}

% ---- Table of contents styling ----------------------------------------------
% By default LaTeX gives dotted leaders to subsections but NOT to sections
% (they print bold with just a right-aligned page number). tocloft lets us turn
% on dots for the section level too, so every TOC line runs "….." to its page.
\usepackage{tocloft}
\renewcommand{\cftsecdotsep}{\cftdotsep}   % dotted leaders for section entries
\renewcommand{\cftsecleader}{\cftdotfill{\cftsecdotsep}}
\renewcommand{\cftsecpagefont}{\normalfont}  % keep page numbers un-bold

% ---- Figures -----------------------------------------------------------------
\usepackage{graphicx}
\graphicspath{{figures/}}

% Landscape (rotated) full-page floats — needed by the ER diagram in §5, which
% is far wider than it is tall and would be unreadable scaled to \textwidth.
\usepackage{rotating}

% ---- Flowcharts / diagrams (TikZ) -------------------------------------------
% Used for the borrow–return workflow diagram (§5). Libraries: geometric shapes
% for decision diamonds, arrows.meta for arrowheads, positioning for relative
% node placement, calc/fit/backgrounds for routing loops and grouping.
\usepackage{tikz}
\usetikzlibrary{shapes.geometric, arrows.meta, positioning, calc, fit, backgrounds}

% ---- Links & bookmarks -------------------------------------------------------
\usepackage{hyperref}
\hypersetup{
  colorlinks=true, linkcolor=black, citecolor=linkblue, urlcolor=linkblue,
  pdflang={th},
}

% Thai labels for auto-generated headings.
\renewcommand{\contentsname}{สารบัญ (Table of Contents)}
\renewcommand{\tablename}{ตาราง}
\renewcommand{\figurename}{รูปที่}

:
%    1) \hypersetup{pdftitle={...}}  — ชื่อเอกสารใน metadata ของ PDF
%    2) \begin{document}
%  (pdftitle ถูกย้ายออกจากไฟล์นี้ เพราะเป็นค่าเฉพาะของแต่ละเอกสาร)
%
%  แก้ที่นี่ที่เดียว = มีผลกับทุกเอกสาร ห้ามก๊อป preamble ไปวางซ้ำ
% =============================================================================
% =============================================================================
%  main.tex  —  เอกสารข้อเสนอโครงการ (Academic Project Proposal)
%  ระบบบริหารจัดการการยืม–คืนอุปกรณ์มหาวิทยาลัย (ULMs)
%
%  IMPORTANT: This document contains Thai text and MUST be compiled with
%  XeLaTeX (not pdflatex). Thai needs a Unicode font + ICU line-breaking.
%      Build:  latexmk main.tex        (uses .latexmkrc -> xelatex)
%      or:     xelatex main.tex        (run twice for the table of contents)
% =============================================================================

\documentclass[11pt, a4paper]{article}

% ---- Fonts & Thai support ----------------------------------------------------
% fontspec lets XeLaTeX use any system font. Laksaman is a complete serif face
% that covers BOTH Thai and Latin glyphs, so one \setmainfont handles the whole
% (Thai + English) document without font switching.
\usepackage{fontspec}

% Thai has no spaces between words, so TeX cannot find line-break points on its
% own. These two XeTeX primitives switch on ICU's dictionary-based Thai word
% segmentation, which is what makes Thai paragraphs wrap correctly.
\XeTeXlinebreaklocale "th"
\XeTeXlinebreakskip = 0pt plus 0.1pt

\setmainfont{Laksaman}

% ---- Page geometry & paragraph style ----------------------------------------
\usepackage[margin=2.5cm]{geometry}
\usepackage{parskip}          % block paragraphs (no indent, gap between) — common
                              % and more readable for Thai documents
\linespread{1.15}             % a little extra leading; Thai has tall tone marks

% ---- Colours (monochrome / grayscale palette) ------------------------------
\usepackage[table]{xcolor}    % 'table' option enables \rowcolor for table headers
% Monochrome (grayscale) palette. Names kept as-is so every reference still
% works; only the tones changed from the original green/blue accents.
\definecolor{kugreen}{HTML}{1A1A1A}   % near-black — section headings & title
\definecolor{kudark}{HTML}{333333}    % dark gray — subsection/subsubsection titles
\definecolor{headbg}{HTML}{EAEAEA}    % light gray background for table header rows
\definecolor{linkblue}{HTML}{404040}  % medium-dark gray — links / citations

% ---- Section heading styling -------------------------------------------------
% 'nobottomtitles*' stops a heading from being printed near the bottom of a
% page: if less than \bottomtitlespace of room is left, the heading (and its
% following text) moves to the next page. This is titlesec's built-in, robust
% replacement for the earlier \needspace hack — the latter injected vertical
% glue that clashed with longtable page-splitting ("Infinite glue shrinkage").
\usepackage[nobottomtitles*]{titlesec}
\titleformat{\section}
  {\color{kugreen}\Large\bfseries}{\thesection}{0.6em}{}
\titleformat{\subsection}
  {\color{kudark}\large\bfseries}{\thesubsection}{0.5em}{}
\titleformat{\subsubsection}
  {\color{kudark}\normalsize\bfseries}{\thesubsubsection}{0.5em}{}
\titlespacing*{\section}{0pt}{1.4\baselineskip}{0.5\baselineskip}

% Minimum room a heading needs at the bottom of a page; if only ~1–3 lines are
% left (less than this), the heading is bumped to the next page instead of being
% stranded above its table/text.
\renewcommand{\bottomtitlespace}{4\baselineskip}

% ---- Hierarchical indentation by heading level ------------------------------
% Each heading level shifts BOTH the heading and the body text under it further
% to the right: section = 0, subsection (x.y) = 1 step, subsubsection (x.y.z) =
% 2 steps. We do this by having the sectioning commands set \leftskip (a global
% left margin that carries through every following line until the next heading
% resets it). \par first, so the *previous* paragraph is not re-indented.
% \leftskip indents plain paragraphs; \@totalleftmargin makes list environments
% (itemize/enumerate) start from the same indent instead of the page margin.
\makeatletter
\newlength{\lvlindent}
\setlength{\lvlindent}{1.8em}          % one indent step
\let\ulmsSecIndent\section
\renewcommand{\section}{\par\global\leftskip=0pt\global\@totalleftmargin=0pt\relax\ulmsSecIndent}
\let\ulmsSubsecIndent\subsection
\renewcommand{\subsection}{\par\global\leftskip=\lvlindent\global\@totalleftmargin=\lvlindent\relax\ulmsSubsecIndent}
\let\ulmsSubsubsecIndent\subsubsection
\renewcommand{\subsubsection}{\par\global\leftskip=2\lvlindent\global\@totalleftmargin=2\lvlindent\relax\ulmsSubsubsecIndent}
\makeatother

% Indent the headings themselves to match (titlesec ignores \leftskip for the
% heading line, so its left margin is set here via the first \titlespacing arg).
\titlespacing*{\subsection}   {\lvlindent} {1.0\baselineskip}{0.4\baselineskip}
\titlespacing*{\subsubsection}{2\lvlindent}{0.8\baselineskip}{0.3\baselineskip}

% ---- Lists (tighter, custom bullets) ----------------------------------------
\usepackage{enumitem}
\setlist[itemize]{leftmargin=1.6em, itemsep=2pt, topsep=3pt}
\setlist[enumerate]{leftmargin=1.9em, itemsep=2pt, topsep=3pt}

% ---- Tables ------------------------------------------------------------------
\usepackage{array}
\usepackage{booktabs}         % \toprule / \midrule / \bottomrule
\usepackage{tabularx}         % auto-width columns that wrap long Thai text
\usepackage{longtable}        % tables that can break across pages
\usepackage{multirow}         % merged cells for the team-structure table

% Left-aligned wrapping column types:
%   L{width} — fixed-width ragged-right paragraph column
%   Y        — flexible ragged-right X column (for tabularx)
\newcolumntype{L}[1]{>{\raggedright\arraybackslash}p{#1}}
\newcolumntype{Y}{>{\raggedright\arraybackslash}X}
\renewcommand{\arraystretch}{1.3}

% Helpers for a styled header row inside tables.
%   \hrow  — starts a header row and shades it (must be FIRST in the row;
%            colortbl allows only one \rowcolor per row).
%   \thead — bold header-cell text.
\newcommand{\hrow}{\rowcolor{headbg}}
\newcommand{\thead}[1]{\textbf{#1}}

% \begin{keeptable} … \end{keeptable} — keeps a table and the bold label above
% it on the SAME page (used by the data dictionary in §5). A tabularx cannot
% split across pages, so LaTeX is free to break between the label and its
% table, stranding the label alone at the bottom of a page; wrapping the pair
% in a minipage makes them one unbreakable block.
%   * \leftskip is zeroed inside: the minipage box is already placed at the
%     current heading indent, so leaving it set would indent the content twice
%     and push the table past the right margin.
%   * \parskip is restored because minipage's \@parboxrestore zeroes it, which
%     would glue the table straight onto the label line.
% Tables inside must use \linewidth (= the minipage width), not \textwidth.
\newenvironment{keeptable}
  {\par\medskip\noindent
   \begin{minipage}{\dimexpr\textwidth-\leftskip\relax}%
   \leftskip=0pt\relax
   \setlength{\parskip}{0.5\baselineskip plus 2pt}}
  {\end{minipage}\par\medskip}

% ---- Table of contents styling ----------------------------------------------
% By default LaTeX gives dotted leaders to subsections but NOT to sections
% (they print bold with just a right-aligned page number). tocloft lets us turn
% on dots for the section level too, so every TOC line runs "….." to its page.
\usepackage{tocloft}
\renewcommand{\cftsecdotsep}{\cftdotsep}   % dotted leaders for section entries
\renewcommand{\cftsecleader}{\cftdotfill{\cftsecdotsep}}
\renewcommand{\cftsecpagefont}{\normalfont}  % keep page numbers un-bold

% ---- Figures -----------------------------------------------------------------
\usepackage{graphicx}
\graphicspath{{figures/}}

% Landscape (rotated) full-page floats — needed by the ER diagram in §5, which
% is far wider than it is tall and would be unreadable scaled to \textwidth.
\usepackage{rotating}

% ---- Flowcharts / diagrams (TikZ) -------------------------------------------
% Used for the borrow–return workflow diagram (§5). Libraries: geometric shapes
% for decision diamonds, arrows.meta for arrowheads, positioning for relative
% node placement, calc/fit/backgrounds for routing loops and grouping.
\usepackage{tikz}
\usetikzlibrary{shapes.geometric, arrows.meta, positioning, calc, fit, backgrounds}

% ---- Links & bookmarks -------------------------------------------------------
\usepackage{hyperref}
\hypersetup{
  colorlinks=true, linkcolor=black, citecolor=linkblue, urlcolor=linkblue,
  pdflang={th},
}

% Thai labels for auto-generated headings.
\renewcommand{\contentsname}{สารบัญ (Table of Contents)}
\renewcommand{\tablename}{ตาราง}
\renewcommand{\figurename}{รูปที่}



\hypersetup{pdftitle={รายงานความก้าวหน้าโครงการ ULMs ครั้งที่ 1}}

% ---- ย่อหน้า ------------------------------------------------------------------
% preamble โหลด parskip ไว้ ซึ่งบังคับ \parindent=0 แล้วเว้นบรรทัดคั่นแทน
% รายงานฉบับนี้ต้องการย่อหน้าแบบไทย คือเยื้องบรรทัดแรก จึงตั้ง \parindent คืน
% และลดช่องไฟระหว่างย่อหน้าลง เพราะถ้าใช้ทั้งเยื้องและเว้นเต็มช่วงพร้อมกัน
% ย่อหน้าจะดูห่างเกินจริง
\setlength{\parindent}{1.5em}
\setlength{\parskip}{0.35\baselineskip plus 2pt}

% preamble ประกาศ \titlespacing* แบบมีดอกจัน ซึ่ง titlesec ตีความว่าห้ามเยื้อง
% ย่อหน้าแรกที่ตามหลังหัวข้อ (ธรรมเนียมฝั่งอังกฤษ) เอกสารไทยเยื้องทุกย่อหน้า
% จึงประกาศซ้ำแบบไม่มีดอกจันด้วยค่าเดิมทุกตัว เพื่อไม่ต้องแก้ preamble ที่
% เอกสารฉบับอื่นใช้ร่วมกัน — ค่าที่ใช้ต้องตรงกับ ulms-preamble.tex เสมอ
\titlespacing{\section}      {0pt}        {1.4\baselineskip}{0.5\baselineskip}
\titlespacing{\subsection}   {\lvlindent} {1.0\baselineskip}{0.4\baselineskip}
\titlespacing{\subsubsection}{2\lvlindent}{0.8\baselineskip}{0.3\baselineskip}

% preamble ตั้ง topsep ของรายการไว้แค่ 3pt ซึ่งแน่นเกินไปเมื่อย่อหน้าเยื้อง
% บรรทัดแรกของย่อหน้าที่ตามหลัง \end{itemize} จะไปเริ่มตรงกับข้อความในหัวข้อย่อย
% พอดี ทำให้อ่านเหมือนเป็นบรรทัดต่อของหัวข้อย่อยสุดท้าย จึงเปิดช่องไฟรอบรายการ
\setlist[itemize]{leftmargin=1.6em, itemsep=2pt, topsep=0.55\baselineskip}
\setlist[enumerate]{leftmargin=1.9em, itemsep=2pt, topsep=0.55\baselineskip}

% ---- การจัดบรรทัด ------------------------------------------------------------
% ภาษาไทยไม่มีการตัดคำด้วยยัติภังค์ และรายงานมีชื่อไฟล์ใน \texttt{} จำนวนมาก
% ซึ่ง TeX ตัดกลางคำไม่ได้ \emergencystretch สั่งให้ยอมยืดช่องไฟแทนการปล่อยล้น
\setlength{\emergencystretch}{3em}

% \zb — จุดตัดบรรทัดกว้างศูนย์ สำหรับแทรกในชื่อยาว ๆ ที่ \texttt ตัดเองไม่ได้
\newcommand{\zb}{\hspace{0pt}}

% ---- การวางรูป --------------------------------------------------------------
% กราฟทั้งสองรูปต้องอยู่ตรงตำแหน่งที่เขียนไว้เป๊ะ ๆ ห้ามลอยข้ามหัวข้อไป
% ถ้าใช้ [!ht] ตามปกติ รูปที่ไม่พอที่จะลอยไปหน้าถัดไป แล้วหัวข้อถัดไปจะขึ้นมา
% แทนที่ กลายเป็นหัวข้อค้างท้ายหน้าโดยไม่มีเนื้อหา — float ให้ตัวเลือก [H]
% ซึ่งบังคับวางตรงนั้นจริง ๆ
\usepackage{float}

% ---- ตัวแปรของรอบรายงาน ------------------------------------------------------
\newcommand{\reportdate}{27 สิงหาคม 2569}
\newcommand{\reportperiod}{24 กรกฎาคม – 27 สิงหาคม 2569}
\newcommand{\actualpct}{52\%}
\newcommand{\planpct}{65\%}

\begin{document}

% =============================================================================
%  TITLE PAGE
% =============================================================================
\begin{titlepage}
\centering
\vspace*{2.8cm}

{\Large\bfseries\color{kugreen} รายงานความก้าวหน้าโครงการ ครั้งที่ 1\par}
\vspace{0.6cm}
{\LARGE\bfseries\color{kugreen} ระบบบริหารจัดการการยืม–คืนอุปกรณ์มหาวิทยาลัย\par}
\vspace{0.35cm}
{\Large\color{kudark} University Lending Management System (ULMs)\par}

\vspace{1.8cm}
\rule{0.72\textwidth}{0.8pt}
\vspace{0.9cm}

\begin{tabular}{@{}r@{\hspace{1.2em}}l@{}}
\textbf{รอบรายงาน}       & \reportperiod \\[0.4em]
\textbf{วันที่รายงาน}     & \reportdate \\[0.4em]
\textbf{ขอบเขตการนับ}    & งานพัฒนาระบบ Sprint 0–5 (ไม่รวมงานก่อน 24 ก.ค.) \\[0.4em]
\textbf{สถานะ Sprint}    & สิ้นสุด Sprint 2 (Sprint ที่ 3 จาก 6) \\[0.4em]
\textbf{ความก้าวหน้า}    & \actualpct\ เทียบแผน \planpct\ \ (ต่ำกว่าแผน 13\%) \\[0.4em]
\textbf{งบประมาณ}        & ใช้ไป 0 บาท จากวงเงินประเมิน 5{,}000 บาท \\[0.4em]
\textbf{กำหนดเสร็จ}      & 1 ตุลาคม 2569 (ยังไม่มีการเลื่อน) \\
\end{tabular}

\vspace{0.9cm}
\rule{0.72\textwidth}{0.8pt}

\vfill
{\large\color{kudark} คณะทำงาน ``miro ปั่น 14 บาท''\par}
\vspace{0.3cm}
{\color{kudark} ภาควิชาวิศวกรรมคอมพิวเตอร์ คณะวิศวกรรมศาสตร์\par}
{\color{kudark} มหาวิทยาลัยเกษตรศาสตร์\par}
\vspace{1.2cm}
\end{titlepage}

% =============================================================================
\section{บทสรุป}

วันรายงานรอบนี้ตรงกับวันสิ้นสุด Sprint 2 ครึ่งแรกของ Sprint ตรงกับช่วงสอบกลางภาคและไม่มีงานเข้าระบบจัดเก็บโค้ดเลย คณะทำงานกลับมาทำงานเต็มกำลังตั้งแต่วันที่ 23 สิงหาคม 2569 และมีงานเข้าทุกวันหลังจากนั้น

สิ่งที่รวมเข้าสายหลักแล้วในรอบนี้มีสามอย่าง คือระบบผู้ดูแลระบบฝั่ง Back-End ครบทั้งชุด สคริปต์ข้อมูลตั้งต้นซึ่งเป็นงานที่ค้างมาตั้งแต่ Sprint 0 และการปรับรุ่นไลบรารีตรวจสอบข้อมูลของทั้งสองฝั่งให้ตรงกัน ซึ่งเป็นเงื่อนไขที่ปิดกั้นการรับสัญญา API มาตั้งแต่ Sprint 1 ผลรวมของทั้งสามอย่างคือหน้าเว็บอ่านข้อมูลผู้ใช้จากฐานข้อมูลจริงได้แล้วเป็นครั้งแรกของโครงการ

งานก้อนใหญ่ที่สุดของรอบนี้คือวงจรยืม–คืนฝั่งเจ้าหน้าที่ กับรายการอุปกรณ์และการค้นหาฝั่งผู้ยืม ทั้งสองส่วนเขียนเสร็จและผ่านการตรวจแล้ว แต่ยังอยู่บนสายงานแยก รายงานฉบับนี้นับเฉพาะสิ่งที่อยู่บนสายหลัก งานที่รอรวมจึงยังไม่ถูกนับ และเป็นเหตุผลหลักที่ผลต่างระหว่างแผนกับผลจริงยังแคบลงเพียงเล็กน้อย

เกณฑ์ ``เสร็จ'' ของ Sprint 2 คือยืมอุปกรณ์ได้ครบวงจรในเส้นทางปกติ และไม่เหลือข้อมูลจำลองในเส้นทางหลัก รอบนี้ไม่ผ่านทั้งสองข้อ เพราะฝั่งผู้ยืมยังไม่มี Procedure สำหรับสร้างคำขอยืม วงจรจึงยังปิดไม่ครบ และหน้าจอทุกหน้ายกเว้นเส้นทางเข้าสู่ระบบยังอ่านข้อมูลจำลอง งานที่เหลือทั้งหมดยกไปเป็นงานลำดับแรกของ Sprint 3

\subsection*{จุดตั้งต้นของการนับความก้าวหน้า}

ตัวเลขในรายงานนี้เริ่มนับจากวันที่ 24 กรกฎาคม 2569 ซึ่งเป็นวันเริ่ม Sprint 0 และเป็นวันเริ่มลงมือพัฒนาระบบ งานที่เสร็จก่อนหน้านั้นถือเป็นงานตั้งต้น และไม่ถูกนับรวมใน \actualpct\ ได้แก่

\begin{itemize}
  \item ข้อเสนอโครงการ — หลักการและเหตุผล · วัตถุประสงค์ · เป้าหมาย · ขอบเขต · แผนการดำเนินงาน · โครงสร้างคณะทำงาน · งบประมาณ
  \item การออกแบบความสามารถของระบบ — ขอบเขตแยกตามบทบาทผู้ใช้ · แผนภาพ Use Case · ผังกระบวนการยืม–คืน · กฎการแบ่งระดับอุปกรณ์ T0–T3 และระบบเครดิต
  \item การออกแบบระบบพื้นฐาน — แผนภาพความสัมพันธ์ของข้อมูลฉบับร่าง และภาพจำลองหน้าจอชุดแรก
\end{itemize}

% =============================================================================
\clearpage
\section{แผนภาพรวมทั้ง 6 Sprint}
\label{sec:overview}

โครงการแบ่งเป็น 6 Sprint บนกรอบเวลา 10 สัปดาห์ ตารางนี้เป็นภาพรวมของทั้งโครงการ รายละเอียดของ Sprint ที่ผ่านมาอยู่ในหัวข้อ \ref{sec:done} และของ Sprint ที่เหลืออยู่ในหัวข้อ \ref{sec:next}

\noindent\begin{tabularx}{\dimexpr\textwidth-\leftskip\relax}{@{} L{2.9cm} Y L{1.5cm} @{}}
\toprule
\hrow \thead{Sprint} & \thead{ธีมและงานหลัก} & \thead{สถานะ} \\
\midrule
\textbf{Sprint 0} \newline 24 – 30 ก.ค. & วางฐาน — ปิดขอบเขต · ติดตั้งโครงสร้างพื้นฐานทั้งสองฝั่ง · แบบจำลองฐานข้อมูลแกนกลาง · ระบบออกแบบและข้อมูลจำลอง & ผ่าน \\
\textbf{Sprint 1} \newline 31 ก.ค. – 13 ส.ค. & ล็อกสัญญาและเปิดทาง — ล็อกสัญญา API · ระบบเข้าสู่ระบบและออกจากระบบ · ระบบตรวจสอบอัตโนมัติ · เกณฑ์การยอมรับของวงจรหลัก & ผ่าน \\
\textbf{Sprint 2} \newline 14 – 27 ส.ค. & วงจรยืม–คืนของอุปกรณ์เคลื่อนย้ายได้ — รายการอุปกรณ์และจำนวนคงเหลือ · สร้างคำขอ อนุมัติ ส่งมอบ รับคืน · ข้อมูลตั้งต้น · เลิกใช้ข้อมูลจำลองในเส้นทางหลัก & ไม่ผ่านเกณฑ์ \newline (ครึ่งแรกตรงกับสอบกลางภาค ยกงานที่เหลือไป Sprint 3) \\
\textbf{Sprint 3} \newline 28 ส.ค. – 10 ก.ย. & กฎธุรกิจรายระดับและการจองสถานที่ (T3) — อนุมัติตามระดับอุปกรณ์ · ผูกหมายเลขชิ้น · สล็อตเวลาและการจองห้อง · ตรวจสภาพและระบบเครดิต \newline รับงานที่ยกมาจาก Sprint 2 เป็นลำดับแรก & เริ่ม 28 ส.ค. \\
\textbf{Sprint 4} \newline 11 – 24 ก.ย. & ข้ามหน่วยงานและงานขัดเงา — ตะกร้าข้ามหน่วยงาน · จองล่วงหน้า · ต่ออายุ · หน้าจอที่เหลือ · การค้นหาและการรองรับมือถือ & ตามแผน \\
\textbf{Sprint 5} \newline 25 ก.ย. – 1 ต.ค. & หยุดเพิ่มความสามารถ — ทดสอบถดถอย · นำขึ้นใช้งานจริง · SRS/SDS และคู่มือ · ซ้อมนำเสนอ & ตามแผน \\
\bottomrule
\end{tabularx}


% =============================================================================
\section{สถานะโครงการโดยรวม (เงิน · งาน · เวลา)}

\noindent\begin{tabularx}{\dimexpr\textwidth-\leftskip\relax}{@{} L{1.5cm} L{3.0cm} L{2.4cm} Y @{}}
\toprule
\hrow \thead{ด้าน} & \thead{แผน} & \thead{ผลจริง} & \thead{การประเมิน} \\
\midrule
เวลา & ผ่านไป 4.9 จาก 10 สัปดาห์ (49\%) & ผ่านไป 4.9 จาก 10 สัปดาห์ (49\%) & ตามปฏิทิน วันเสร็จยังเป็น 1 ตุลาคม 2569 \\
งาน & \planpct & \actualpct & ต่ำกว่าแผน 13\% จากช่วงสอบกลางภาคที่ไม่มีงานเข้าครึ่ง Sprint และจากงานที่เสร็จแล้วแต่ยังไม่รวมเข้าสายหลัก รายละเอียดรายมิติอยู่ในหัวข้อ \ref{sec:dimension} \\
เงิน & 5{,}000 บาท (วงเงินประเมินทั้งโครงการ) & 0 บาท (0\%) & ยังไม่มีรายจ่าย รายการที่ประเมินไว้คือค่าเช่า Cloud/Hosting และฐานข้อมูลกลาง ซึ่งจะเกิดขึ้นเมื่อนำระบบขึ้นใช้งานจริงใน Sprint 5 \\
\bottomrule
\end{tabularx}

% =============================================================================
\section{ความก้าวหน้าแยกตามมิติของงาน}
\label{sec:dimension}

การแบ่งมิติเป็นไปตามสไลด์ Week 06 คือแยกส่วนออกแบบ ส่วนพัฒนา ส่วนทดสอบ และส่วนเอกสาร ค่าในคอลัมน์ ``จริง'' ประเมินจากผลงานที่ตรวจสอบได้ในระบบจัดเก็บโค้ด

วันรายงานรอบนี้ตรงกับวันสิ้นสุด Sprint 2 พอดี ค่าในคอลัมน์ ``แผน'' จึงอ่านได้ตรงจากแถว Sprint 2 ในตารางภาคผนวก ไม่ต้องสอดแทรกค่า ส่วนค่าในคอลัมน์ ``จริง'' นับเฉพาะงานที่รวมเข้าสายหลักแล้ว งานที่เขียนเสร็จแต่ยังรอรวมไม่ถูกนับในตารางนี้ และระบุไว้ในคอลัมน์ข้อเท็จจริงประกอบแทน

\noindent\begin{tabularx}{\dimexpr\textwidth-\leftskip\relax}{@{} L{3.2cm} L{1.1cm} L{1.0cm} L{1.0cm} L{1.2cm} Y @{}}
\toprule
\hrow \thead{มิติงาน} & \thead{น้ำหนัก} & \thead{แผน} & \thead{จริง} & \thead{ผลต่าง} & \thead{ข้อเท็จจริงประกอบ} \\
\midrule
Back-End Design \newline (API + Database Model) & 10\% & 95\% & 86\% & $-9$ & สัญญาของกลุ่มยืนยันตัวตนและกลุ่มผู้ดูแลระบบอยู่บนสายหลักแล้ว สัญญาของวงจรยืม–คืนเขียนเสร็จแต่ยังรอรวม \\
Front-End Design \newline (UI/UX + Business Process) & 10\% & 85\% & 80\% & $-5$ & ออกแบบหน้าจอที่มีเนื้อหาจริงแล้ว 14 จาก 30 หน้า \\
Back-End Dev & 30\% & 60\% & 42\% & $-18$ & ระบบยืนยันตัวตนและระบบผู้ดูแลระบบใช้งานได้จริง ยังไม่มีวงจรยืม–คืนบนสายหลัก \\
Front-End Dev & 25\% & 65\% & 54\% & $-11$ & เส้นทางเข้าสู่ระบบอ่านข้อมูลจริงแล้ว อีก 13 หน้ายังอ่านข้อมูลจำลอง \\
Test (Front-End + Back-End) & 15\% & 35\% & 14\% & $-21$ & ชุดทดสอบฝั่ง Back-End 55 กรณีบนสายหลัก แต่ระบบตรวจสอบอัตโนมัติยังไม่ได้เรียกใช้ และยังไม่มีตัวรันฝั่ง Front-End \\
Documentation & 10\% & 70\% & 68\% & $-2$ & ยังไม่มี SRS และ SDS ฉบับทางการ \\
\midrule
\hrow \thead{รวมถ่วงน้ำหนัก} & \thead{100\%} & \thead{\planpct} & \thead{\actualpct} & \thead{$-13$} & \\
\bottomrule
\end{tabularx}

มิติที่ต้องเฝ้าระวังมีสองมิติ มิติแรกคือ Back-End Dev ซึ่งมีน้ำหนักสูงสุดที่ 30\% ขยับขึ้นมากที่สุดในบรรดาทุกมิติมาอยู่ที่ 42\% แต่ยังห่างจากแผนอยู่ $-18$ ข้อสังเกตคืองานที่ทำให้มิตินี้เดินต่อได้เขียนเสร็จหมดแล้วและรออยู่บนสายงานแยก การรวมงานเหล่านั้นเข้าสายหลักจึงเป็นงานที่ให้ผลตอบแทนสูงที่สุดในตอนนี้

มิติที่สองคือ Test ซึ่งห่างจากแผนมากที่สุดในรอบนี้ที่ $-21$ สาเหตุไม่ใช่ว่าไม่มีชุดทดสอบ แต่เป็นเพราะระบบตรวจสอบอัตโนมัติยังไม่ได้เรียกชุดทดสอบที่มีอยู่ให้ทำงาน และฝั่ง Front-End ยังไม่มีตัวรันชุดทดสอบเลย ทั้งสองอย่างใช้เวลาไม่มาก แต่ต้องทำก่อนที่ปริมาณโค้ดจะโตกว่านี้

\begin{figure}[H]
\centering
\begin{tikzpicture}[x=0.062cm, y=1cm]
  % ---- กริดแนวตั้งและสเกลแกนนอน ----
  \foreach \x in {0,20,40,60,80,100} {
    \draw[gray!25] (\x,-0.30) -- (\x,6.15);
    \node[below, font=\scriptsize] at (\x,-0.30) {\x};
  }
  \node[below, font=\scriptsize] at (50,-0.68) {\% ความก้าวหน้าของมิตินั้น};
  \draw[thick] (0,-0.30) -- (0,6.15);

  % ---- แท่งรายมิติ: y / ชื่อ / น้ำหนัก / แผน / จริง ----
  \foreach \y/\name/\wt/\plan/\act in {%
      5.6/{Back-End Design}/10/95/86,
      4.6/{Front-End Design}/10/85/80,
      3.6/{Back-End Dev}/30/60/42,
      2.6/{Front-End Dev}/25/65/54,
      1.6/{Test}/15/35/14,
      0.6/{Documentation}/10/70/68} {
    \node[left, xshift=-0.16cm, font=\scriptsize, align=right] at (0,\y)
      {\name\\[-3pt]{\tiny น้ำหนัก \wt\%}};
    \fill[gray!45]  (0,{\y+0.05}) rectangle (\plan,{\y+0.30});
    \fill[kugreen]  (0,{\y-0.30}) rectangle (\act,{\y-0.05});
    \node[right, font=\tiny, gray!45!black] at (\plan,{\y+0.175}) {\plan};
    \node[right, font=\tiny, kugreen]       at (\act,{\y-0.175}) {\act};
  }

  % ---- คำอธิบายสัญลักษณ์ ----
  \fill[gray!45] (0,-1.36) rectangle (7,-1.54);
  \node[right, font=\scriptsize] at (9,-1.45) {แผน};
  \fill[kugreen] (0,-1.76) rectangle (7,-1.94);
  \node[right, font=\scriptsize] at (9,-1.85) {ผลจริง};
\end{tikzpicture}
\vspace{0.6em}
\caption{ความก้าวหน้ารายมิติ ณ \reportdate\ เทียบกับค่าเป้าหมายตามแผน
  ตัวเลขใต้ชื่อมิติคือน้ำหนักที่ใช้ถ่วงเมื่อรวมเป็นค่าเดียว
  แท่งผลจริงนับเฉพาะงานที่รวมเข้าสายหลักแล้ว
  มิติที่ห่างจากแผนมากที่สุดคือ Test ที่ $-21$ รองลงมาคือ Back-End Dev ที่ $-18$
  ซึ่งเป็นมิติที่มีน้ำหนักสูงสุดในโครงการ}
\label{fig:dims}
\end{figure}

% =============================================================================
\section{กราฟความก้าวหน้าเทียบแผน}

เส้นแผนสร้างจากแผน Sprint โดยตรง ไม่ได้กำหนดค่าด้วยมือ คือกำหนดค่าเป้าหมายรายมิติ ณ วันสิ้นสุดของแต่ละ Sprint จากรายการงานที่ Sprint นั้นรับไว้ แล้วคูณน้ำหนักรวมกันเป็นค่าสะสม ตารางค่าทั้งหมดอยู่ในภาคผนวก

แกนนอนเริ่มที่สัปดาห์ที่ 0 ซึ่งตรงกับ 24 กรกฎาคม 2569 เป็นวันเริ่มลงมือพัฒนา ไม่ใช่วันเริ่มโครงการ งานข้อเสนอและงานออกแบบที่เสร็จก่อนหน้านั้นจึงไม่ปรากฏบนกราฟ และ Documentation ที่เริ่มด้วยค่า 10\% 

จุดบนเส้นแผนวางที่ขอบ Sprint ไม่ใช่ตามสัปดาห์ปฏิทิน เพราะคณะทำงานประเมินเกณฑ์ ``เสร็จ'' ที่ขอบ Sprint ความชันของแต่ละช่วงจึงสะท้อนปริมาณงานที่ Sprint นั้นรับไว้จริง วันรายงานรอบนี้ตรงกับขอบ Sprint 2 พอดี ค่าแผนจึงอ่านจากจุดบนเส้นแผนได้ตรง ๆ

รูปร่างของเส้นผลจริงในรอบนี้บอกเรื่องสองเรื่อง ช่วงราบระหว่างวันที่ 13 ถึง 23 สิงหาคม 2569 คือช่วงสอบกลางภาคซึ่งไม่มีงานเข้าระบบจัดเก็บโค้ดเลย และช่วงที่ไต่ขึ้นหลังจากนั้นคือห้าวันสุดท้ายของ Sprint ที่คณะทำงานกลับมาเต็มกำลัง ความชันของช่วงหลังนี้ชันกว่าเส้นแผนในช่วงเดียวกัน แต่ระยะเวลาสั้นเกินกว่าจะไล่ส่วนที่หายไปได้ทัน แท่งแนวตั้งที่วันรายงานคือระยะห่างระหว่างแผนกับผลจริงในรอบนี้

\begin{figure}[H]
\centering
\begin{tikzpicture}[x=1.20cm, y=0.058cm]
  % ---- กริดแนวนอน ----
  \foreach \y in {0,20,40,60,80,100}
    \draw[gray!22] (0,\y) -- (10.2,\y);

  % ---- เส้นประกำกับขอบ Sprint ----
  \foreach \x in {0.86, 2.86, 4.86, 6.86, 8.86}
    \draw[gray!40, dashed, thin] (\x,0) -- (\x,108);

  % ---- ป้ายชื่อ Sprint ด้านบน ----
  \foreach \x/\lbl in {0.43/{S0}, 1.86/{S1}, 3.86/{S2}, 5.86/{S3}, 7.86/{S4}, 9.36/{S5}}
    \node[font=\scriptsize\bfseries, gray!50!black] at (\x,114) {\lbl};

  % ---- แกน ----
  \draw[->, thick] (0,0) -- (10.6,0);
  \draw[->, thick] (0,0) -- (0,122) node[above, font=\scriptsize] {\% ความก้าวหน้าสะสม};
  \foreach \y in {0,20,40,60,80,100}
    \draw (0,\y) -- (-0.10,\y) node[left, font=\scriptsize] {\y};

  % ---- ขีดและวันที่ตามขอบ Sprint ----
  \foreach \x/\lbl in {0/{24 ก.ค.}, 0.86/{30 ก.ค.}, 2.86/{13 ส.ค.},
                       4.86/{27 ส.ค.}, 6.86/{10 ก.ย.}, 8.86/{24 ก.ย.}, 9.86/{1 ต.ค.}}
    {\draw (\x,0) -- (\x,-3);
     \node[rotate=45, anchor=north east, font=\scriptsize] at (\x,-4) {\lbl};}

  % ---- เส้นแนวตั้งของวันรายงาน (27 ส.ค. = สัปดาห์ที่ 4.86 = ขอบ Sprint 2) ----
  \draw[gray!55, dotted, thin] (4.86,0) -- (4.86,108);

  % ---- เส้นแผน — จุดวางที่ขอบ Sprint ----
  \draw[thick, kudark, densely dashed]
    plot coordinates {(0,1) (0.86,12) (2.86,42) (4.86,65) (6.86,78) (8.86,94) (9.86,100)};
  \foreach \p in {(0,1), (0.86,12), (2.86,42), (4.86,65), (6.86,78), (8.86,94), (9.86,100)}
    \fill[kudark] \p circle (1.6pt);

  % ---- เส้นผลจริง — จุดคือวันที่วัดความก้าวหน้า ----
  % ช่วง 2.86 -> 4.29 ราบ เพราะช่วงสอบกลางภาคไม่มีงานเพิ่ม
  % 4.29 -> 4.86 คือ 23–27 ส.ค. ห้าวันที่คณะทำงานกลับมาเต็มกำลัง
  \draw[very thick, kugreen]
    plot coordinates {(0,0) (2.0,24) (2.43,32) (2.86,43) (4.29,43) (4.86,52)};
  \foreach \p in {(0,0), (2.0,24), (2.43,32), (2.86,43), (4.29,43), (4.86,52)}
    \fill[kugreen] \p circle (2.2pt);

  % ---- ระยะห่างแผน–ผลจริง ณ วันรายงาน ----
  \draw[gray!70, thick, |-|] (4.86,52) -- (4.86,65);
  \node[right, font=\tiny, gray!70!black] at (4.92,58.5) {$-13$};

  % ---- ป้ายจุดปัจจุบัน ----
  % จัดชิดขวาให้จบก่อนถึงจุดวัด และแบ่งสองบรรทัด เพื่อไม่ให้ข้อความล้น
  % ออกไปทับแกนตั้งทางซ้าย
  \node[anchor=south east, font=\scriptsize, gray!55!black, align=right] at (2.80,45)
    {13 ส.ค.\\43\% · แผน 42\%};
  % fill=white so the label masks the Sprint 3 dashed line it now sits across
  \node[anchor=north west, font=\scriptsize, kugreen, align=left,
        fill=white, inner sep=1.5pt] at (4.95,48)
    {27 ส.ค. (วันรายงาน)\\ผลจริง 52\% · แผน 65\%};

  % ---- คำอธิบายสัญลักษณ์ (บังกริดด้วยพื้นขาวก่อนวาด) ----
  \fill[white] (5.85,4) rectangle (10.2,26);
  \draw[thick, kudark, densely dashed] (6.0,19) -- (6.7,19);
  \node[right, font=\scriptsize] at (6.8,19) {แผน (Planned)};
  \draw[very thick, kugreen] (6.0,9) -- (6.7,9);
  \node[right, font=\scriptsize] at (6.8,9) {ผลจริง (Actual)};
\end{tikzpicture}
\caption{กราฟ เปรียบเทียบความก้าวหน้าตามแผนกับผลการดำเนินงานจริง
  แกนนอนวางตามขอบ Sprint ทั้ง 6 ช่วง (เส้นประแนวตั้ง)
  จุดบนเส้นผลจริงคือวันที่วัดความก้าวหน้า
  วันรายงาน 27 สิงหาคม 2569 ตรงกับขอบ Sprint 2 พอดี
  ช่วงราบระหว่าง 13 ถึง 23 สิงหาคม คือช่วงสอบกลางภาคที่ไม่มีงานเพิ่ม
  ช่วงที่ไต่ขึ้นหลังจากนั้นคือห้าวันสุดท้ายของ Sprint}
\label{fig:scurve}
\end{figure}

% =============================================================================
\section{ผลการดำเนินงาน Sprint 0 ถึงปัจจุบัน}
\label{sec:done}

\noindent\begin{tabularx}{\dimexpr\textwidth-\leftskip\relax}{@{} L{1.1cm} Y L{2.9cm} @{}}
\toprule
\hrow \thead{ส่วนงาน} & \thead{งานตามแผน Sprint 0 (24 – 30 ก.ค.)} & \thead{ผล} \\
\midrule
SA & ปิดขอบเขตของโครงการ & เสร็จ \\
BE & ติดตั้ง NestJS + tRPC + Prisma + PostgreSQL & เสร็จ \\
BE & แบบจำลองฐานข้อมูลแกนกลาง พร้อม Migration ขึ้นฐานข้อมูลจริง & เสร็จ \\
BE & สัญญา tRPC ของวงจรหลัก & เลื่อนไป Sprint 1 \\
BE & ข้อมูลตั้งต้น 2 หน่วยงาน & เลื่อนไป Sprint 2 \\
FE & ติดตั้ง Vite + React + TypeScript + Router และโครงยืนยันตัวตน & เสร็จ \\
FE & ระบบออกแบบและ Wireframe หน้าเข้าสู่ระบบกับเปลือกระบบ & เสร็จ \\
FE & ชั้นข้อมูลจำลองสำหรับพัฒนาหน้าจอ & เสร็จ \\
QA & แผนการทดสอบและเกณฑ์การยอมรับ & เลื่อนไป Sprint 1 \\
PM & ตั้งกระดานงาน พิธีการประจำ Sprint และเกณฑ์ ``เสร็จ'' & เสร็จ \\
\bottomrule
\end{tabularx}

\noindent\begin{tabularx}{\dimexpr\textwidth-\leftskip\relax}{@{} L{1.1cm} Y L{3.4cm} @{}}
\toprule
\hrow \thead{ส่วนงาน} & \thead{งานตามแผน Sprint 1 (31 ก.ค. – 13 ส.ค.)} & \thead{ผล} \\
\midrule
SA & ล็อกสัญญา tRPC และคอมมิตไฟล์สัญญาเข้าระบบจัดเก็บโค้ด & เสร็จ รอทีม Front-End รับไปใช้แทนสัญญาที่เขียนด้วยมือ \\
BE & เพิ่มข้อกำหนดที่ยังขาดในสคีมา & บางส่วน ปรับเป็น 27 ตาราง 9 Enum ปิดเรื่องการผูกชิ้นอุปกรณ์แล้ว ส่วนคีย์ไม่ซ้ำบนอีเมลกับรหัสผู้ใช้เขียนเสร็จแล้วแต่ยังรอรวมเข้าสายหลัก \\
BE & ระบบเข้าสู่ระบบและออกจากระบบ & เสร็จ ทดสอบครบวงบนฐานข้อมูลจริงแล้ว และรวมเข้าสายหลักระหว่าง Sprint 2 \\
BE & ข้อมูลตั้งต้น 2 หน่วยงาน & ยังไม่เสร็จ ยกไป Sprint 2 \\
FE & รวมระบบตรวจสอบอัตโนมัติ (CI) เข้าสายหลัก & เสร็จ \\
QA & เกณฑ์การยอมรับของวงจรหลัก & บางส่วน เขียนเป็นโค้ดตรวจได้แล้ว ยังไม่มีเอกสารแผนการทดสอบ \\
\midrule
\hrow \multicolumn{3}{@{}l@{}}{\thead{งานที่ทำเพิ่มนอกแผน Sprint 1 (ดึงงานของ Sprint 2 และ 3 ขึ้นมาทำก่อน)}} \\
FE & หน้าจอผู้ดูแลระบบ 6 หน้า และหน้าจอผู้ยืม 3 หน้า พร้อมหน้าโปรไฟล์ & เสร็จ (ใช้ข้อมูลจำลอง) \\
FE & ชั้นเชื่อมต่อ tRPC ฝั่งไคลเอนต์ และไลบรารีกฎธุรกิจ & เสร็จ \\
\bottomrule
\end{tabularx}

เกณฑ์ ``เสร็จ'' ของ Sprint 1 ผ่าน 3 ข้อจาก 4 ข้อ คือสัญญา API ระบบยืนยันตัวตน และระบบตรวจสอบอัตโนมัติ เหลือข้อมูลตั้งต้นที่ยกไปเป็นงานลำดับแรกของ Sprint 2

ครึ่งแรกของ Sprint 2 คือวันที่ 14 ถึง 22 สิงหาคม 2569 ตรงกับช่วงสอบกลางภาคทั้งช่วง ไม่มีงานใดเข้าระบบจัดเก็บโค้ด ผลของ Sprint 2 ทั้งหมดในตารางถัดไปจึงเกิดขึ้นภายในห้าวันทำการสุดท้ายคือวันที่ 23 ถึง 27 สิงหาคม 2569

\noindent\begin{tabularx}{\dimexpr\textwidth-\leftskip\relax}{@{} L{1.1cm} Y L{3.4cm} @{}}
\toprule
\hrow \thead{ส่วนงาน} & \thead{งานตามแผน Sprint 2 (14 – 27 ส.ค.)} & \thead{ผล} \\
\midrule
BE & ข้อมูลตั้งต้น 2 หน่วยงาน พร้อมผู้ใช้ครบทุกบทบาท ห้อง และอุปกรณ์ตัวอย่าง & บางส่วน ส่วนของบัญชีผู้ใช้และระดับเครดิตอยู่บนสายหลักแล้ว ส่วนของอุปกรณ์และห้องเขียนเสร็จแล้วแต่ยังรอรวม \\
BE & สัญญาและ Procedure ของรายการอุปกรณ์ การค้นหา และจำนวนคงเหลือ & เสร็จ รอรวมเข้าสายหลัก \\
BE & สัญญาและ Procedure ของการสร้างคำขอ อนุมัติ ผูกชิ้น ส่งมอบ รับคืน & บางส่วน ฝั่งเจ้าหน้าที่เสร็จแล้วรอรวม ฝั่งผู้ยืมสำหรับสร้างคำขอยังไม่เริ่ม \\
FE & รับไฟล์สัญญาไปใช้แทนสัญญาที่เขียนด้วยมือ & ยังไม่เสร็จ ปรับรุ่นไลบรารีตรวจสอบข้อมูลให้ตรงกันแล้วซึ่งเป็นเงื่อนไขที่ปิดกั้นอยู่ เหลือขั้นตอนรับไฟล์ \\
FE & เปลี่ยนหน้าจอที่มีอยู่จากข้อมูลจำลองเป็นข้อมูลจริง & บางส่วน เส้นทางเข้าสู่ระบบและข้อมูลผู้ใช้ปัจจุบันใช้ข้อมูลจริงแล้ว หน้าที่เหลือยังใช้ข้อมูลจำลอง \\
QA & ตัวรันชุดทดสอบฝั่ง Front-End และงานตรวจชุดทดสอบในระบบตรวจสอบอัตโนมัติ & ยังไม่เสร็จ ยกไป Sprint 3 \\
\midrule
\hrow \multicolumn{3}{@{}l@{}}{\thead{งานที่ทำเพิ่มนอกแผน Sprint 2}} \\
BE & ระบบผู้ดูแลระบบฝั่ง Back-End ครบชุด — จัดการบัญชีผู้ใช้ ตั้งค่าระบบ สถานะระบบ และบันทึกการตรวจสอบ & เสร็จ อยู่บนสายหลักแล้ว \\
BE & ระบบจัดการคลังอุปกรณ์และห้องของเจ้าหน้าที่ การตรวจสภาพ การซ่อม และการอัปโหลดรูปภาพ & เสร็จ รอรวมเข้าสายหลัก \\
FE & หน้าจอผู้ยืมเพิ่มอีก 5 หน้า รวมหน้ารับของและหน้าใช้ห้อง & เสร็จ โดยยังใช้ข้อมูลจำลอง และ 2 หน้าหลังรอรวมเข้าสายหลัก \\
FE & ปรับรุ่นไลบรารีร่วมของทั้งสองฝั่ง และปิดช่องโหว่ความปลอดภัยที่ตรวจพบในไลบรารีที่ใช้ & เสร็จ \\
\bottomrule
\end{tabularx}

เกณฑ์ ``เสร็จ'' ของ Sprint 2 มีสองข้อ คือยืมอุปกรณ์ได้ครบวงจรในเส้นทางปกติ และไม่เหลือข้อมูลจำลองในเส้นทางหลัก รอบนี้ไม่ผ่านทั้งสองข้อ สาเหตุแยกได้เป็นสองส่วน ส่วนแรกคือเวลาที่หายไปครึ่ง Sprint จากช่วงสอบกลางภาค ส่วนที่สองคือทุกคนทำงานของตัวเองอยู่บนสายงานแยกจนถึงวันสุดท้ายของ Sprint งานที่เสร็จแล้วจึงยังไม่ได้ประกอบเข้าด้วยกัน และยังไม่มีใครพิสูจน์ได้ว่าเมื่อประกอบแล้วจะทำงานร่วมกันได้ คณะทำงานจึงกำหนดให้การรวมงานเข้าสายหลักเป็นงานลำดับแรกของสัปดาห์ถัดไป และปรับให้รวมงานบ่อยขึ้นแทนการรวมครั้งเดียวตอนจบ Sprint

% =============================================================================
\section{ผลงานที่ส่งมอบแล้ว}

หัวข้อนี้รายงานเฉพาะสิ่งที่ทำงานได้จริงและตรวจสอบได้ในระบบจัดเก็บโค้ด ไม่รวมกิจกรรมระหว่างดำเนินงาน

\subsection{ฝั่ง Back-End}

\begin{itemize}
  \item โครงสร้างพื้นฐาน — NestJS 11 ประกอบกับ Prisma 7, \texttt{nestjs-trpc} และตัวเชื่อม PostgreSQL · ตั้งค่า CORS แบบระบุ Origin เจาะจงพร้อมรองรับคุกกี้ยืนยันตัวตน
  \item ฐานข้อมูล — 27 ตาราง 9 Enum พร้อม Migration ที่สร้างตารางได้จริง ครอบคลุมอุปกรณ์ครบทั้งสี่ระดับ โดยตารางห้องและตารางการจองรองรับ T3 อยู่แล้ว ทั้งประเภททรัพยากร ข้อมูลห้อง และช่วงเวลาเริ่ม–สิ้นสุดของการจอง
  \item ระบบยืนยันตัวตนที่ใช้งานได้จริง — เข้าสู่ระบบด้วยอีเมลมหาวิทยาลัยหรือรหัสผู้ใช้ · เก็บรหัสผ่านด้วย scrypt · ออกคุกกี้เซสชันแบบ httpOnly ที่หน้าเว็บอ่านไม่ได้ · ตรวจโทเคนและอ่านบทบาทสดจากฐานข้อมูลทุกคำขอ · ออกจากระบบ
  \item ชั้นควบคุมสิทธิ์ — บังคับเข้าสู่ระบบก่อน และแยกสิทธิ์อีก 3 ระดับคือเจ้าหน้าที่ ผู้อนุมัติ และผู้ดูแลระบบ · แยกกรณียังไม่เข้าสู่ระบบออกจากกรณีเข้าสู่ระบบแล้วแต่สิทธิ์ไม่พอ
  \item ระบบผู้ดูแลระบบครบชุด — ค้นและแสดงบัญชีผู้ใช้แบบแบ่งหน้า · สร้าง แก้ไข เปลี่ยนบทบาท ตั้งรหัสผ่านใหม่ ระงับและปลดระงับบัญชี · อ่านและแก้ค่าตั้งของระบบการยืม · อ่านสถานะระบบและงานตามเวลา · ค้นบันทึกการตรวจสอบ
  \item สคริปต์ข้อมูลตั้งต้น — สร้างระดับเครดิตให้ครอบคลุมคะแนน 0 ถึง 100 โดยไม่มีช่องว่าง และสร้างบัญชีผู้ใช้ครบทุกบทบาท เขียนให้เรียกซ้ำได้โดยไม่พัง จึงใช้กับฐานข้อมูลที่มีข้อมูลอยู่แล้วได้
  \item สัญญา API เป็นไฟล์กลางที่สร้างจากโค้ดจริง — มีคำสั่งสร้างใหม่และคำสั่งตรวจว่าไฟล์ยังตรงกับโค้ด ทำให้สัญญาไม่มีวันคลาดจากสิ่งที่ระบบทำจริง
  \item ผลการทดสอบ — ชุดทดสอบบนสายหลัก 55 กรณี ครอบคลุมการอ่านและตรวจโทเคนเซสชัน การเข้ารหัสรหัสผ่าน การแปลงข้อมูลบัญชีผู้ใช้ และวงจรเข้าสู่ระบบ ในจำนวนนี้ 48 กรณีทำงานได้โดยไม่ต้องต่อฐานข้อมูล อีก 7 กรณีต้องต่อ PostgreSQL จริง
\end{itemize}

ข้อจำกัดที่เหลือคือสายหลักยังไม่มีวงจรยืม–คืนเลย สัญญาปัจจุบันครอบคลุมเฉพาะกลุ่มยืนยันตัวตนกับกลุ่มผู้ดูแลระบบ ข้อมูลตั้งต้นยังมีเฉพาะส่วนที่การเข้าสู่ระบบต้องใช้ ยังไม่มีอุปกรณ์และห้อง คีย์ไม่ซ้ำบนอีเมลและรหัสผู้ใช้ยังไม่ถูกรวมเข้าสายหลัก ฐานข้อมูลจึงยังยอมรับอีเมลซ้ำได้ และระบบตรวจสอบอัตโนมัติยังไม่ได้เรียกชุดทดสอบที่มีอยู่ให้ทำงาน

\subsection{ฝั่ง Front-End}

\begin{itemize}
  \item โครงสร้างพื้นฐาน — Vite + React 18 + TypeScript + React Router · เส้นทางหน้าจอ 30 เส้นทาง พร้อมชั้นคุ้มกันสิทธิ์ตามบทบาท 4 กลุ่ม · โหลดหน้าจอแบบแยกก้อนทุกหน้า
  \item ระบบออกแบบ — Tailwind CSS พร้อมโทเคนสีและส่วนประกอบพื้นฐานที่ใช้ซ้ำได้ รวมตารางข้อมูล กราฟ กล่องโต้ตอบ และแผงเลื่อน · รองรับไทย–อังกฤษ และโหมดสว่าง–มืด
  \item หน้าจอที่มีเนื้อหาจริง 14 จาก 30 หน้า — ผู้ดูแลระบบ 6 หน้า (ผู้ใช้ · บันทึกการตรวจสอบ · แผงสรุป · สถานะระบบ · ตั้งค่าระบบ · รายงาน) · ผู้ยืม 6 หน้า (รายการอุปกรณ์ · รายละเอียดอุปกรณ์ · รายการยืมของฉัน · สร้างคำขอยืม · รายการห้องและจองห้อง ซึ่งเป็นจุดตั้งต้นของงาน T3) · โปรไฟล์ผู้ใช้ · เข้าสู่ระบบ
  \item เส้นทางแรกที่อ่านข้อมูลจริง — หน้าเข้าสู่ระบบส่งรหัสผ่านไปตรวจที่ฐานข้อมูลจริง เมื่อเปิดเว็บระบบจะถามฝั่ง Back-End ว่าผู้ใช้คนปัจจุบันเป็นใครแล้วแสดงชื่อกับบทบาทที่ได้ และปุ่มออกจากระบบลบเซสชันที่ฝั่ง Back-End จริง
  \item ไลบรารีตรวจสอบข้อมูลตรงรุ่นกับฝั่ง Back-End แล้ว — เป็นเงื่อนไขที่ปิดกั้นการรับไฟล์สัญญามาตั้งแต่ Sprint 1 เมื่อตรงรุ่นแล้ว การรับไฟล์สัญญาจึงเหลือเพียงขั้นตอนเปลี่ยนที่มาของชนิดข้อมูล
  \item ไลบรารีกฎธุรกิจ — ครอบคลุมช่วงเวลายืม เวลาตัดคืน 17:00 น. การนับช่องจองพร้อมกันของ T3 การตรวจไฟล์อัปโหลด และข้อความแสดงข้อผิดพลาด พร้อมสคริปต์ตรวจเกณฑ์การยอมรับ
\end{itemize}

ข้อจำกัดที่เหลือคือมีเพียงเส้นทางเข้าสู่ระบบเท่านั้นที่อ่านข้อมูลจริง อีก 13 หน้ายังอ่านข้อมูลจากไฟล์จำลองทั้งหมด ชนิดข้อมูลของสัญญายังมาจากไฟล์ที่ฝั่ง Front-End เขียนขึ้นเอง ไม่ใช่ไฟล์ที่ฝั่ง Back-End สร้าง ยังไม่มีตัวรันชุดทดสอบฝั่ง Front-End และยังเหลือหน้าว่างอีก 16 หน้า ครอบคลุมเจ้าหน้าที่ 7 หน้า ผู้อนุมัติ 2 หน้า รายงาน 2 หน้า ผู้ยืมที่เหลือ 4 หน้า และตั้งค่าผู้ดูแลระบบ 1 หน้า

\subsection{เครื่องมือและกระบวนการ}

ระบบตรวจสอบอัตโนมัติทำงานทั้งบนคำขอรวมโค้ดและการนำโค้ดเข้าสายหลัก ตรวจความสอดคล้องของเวอร์ชันไลบรารีร่วม แล้วตรวจการติดตั้ง ชนิดข้อมูล และการ Build ของทั้งสองฝั่ง ขณะนี้อยู่บนสายหลักแล้วและการตรวจล่าสุดผ่าน

ข้อจำกัดที่เหลือมีสองข้อ ข้อแรกคืองานตรวจความสอดคล้องของเวอร์ชันยังตั้งไว้ให้รายงานผลโดยไม่ขวางการรวมโค้ด ซึ่งตั้งไว้ชั่วคราวตอนที่ทั้งสองฝั่งยังใช้ไลบรารีคนละรุ่นหลัก ตอนนี้ทั้งสองฝั่งตรงรุ่นกันแล้ว จึงเปลี่ยนให้งานนี้ขวางการรวมโค้ดได้ทันทีเพื่อไม่ให้เวอร์ชันเคลื่อนออกจากกันอีก ข้อที่สองคือระบบตรวจสอบอัตโนมัติตรวจแค่ว่าโค้ดคอมไพล์ผ่านและ Build ได้ ยังไม่ได้เรียกชุดทดสอบที่มีอยู่ให้ทำงาน ชุดทดสอบจึงพังได้โดยไม่มีใครรู้

\subsection{เอกสารและการออกแบบ}

\begin{itemize}
  \item ข้อเสนอโครงการฉบับสมบูรณ์ และการออกแบบฐานข้อมูลพร้อมพจนานุกรมข้อมูลครบทุกตาราง
  \item เอกสารประกอบการตัดสินใจเรื่องสัญญา API — คู่มือหลักการ วาระประชุมพร้อมโค้ดตัวอย่าง และแบบฟอร์มลงมติ
  \item รายงานผลการทดลองสถาปัตยกรรม — ระบุจุดที่ต้องแก้เพื่อให้ทั้งสองฝั่งเชื่อมต่อกันได้ 19 จุด เรียงตามความรุนแรง
  \item เอกสารส่งมอบสัญญาให้ทีม Front-End — ระบุสิ่งที่ต้องแก้ ลำดับการแก้ และผลกระทบต่องานที่ทำอยู่
\end{itemize}

ข้อจำกัดที่เหลือคือยังไม่มี SRS ฉบับทางการ ยังไม่มี SDS นอกเหนือจากบทที่ว่าด้วยฐานข้อมูล และยังไม่มีคู่มือผู้ใช้กับคู่มือติดตั้ง ทั้งสามรายการอยู่ในแผน Sprint 5

% =============================================================================
\section{แผนการดำเนินงานถัดไป}
\label{sec:next}

\subsection{แผนสัปดาห์ถัดไป (28 สิงหาคม – 3 กันยายน 2569)}

สองวันแรกไม่มีการเขียนความสามารถใหม่ เป็นการรวมงานที่เสร็จแล้วเข้าสายหลักทั้งหมด เพื่อให้ปัญหาจากการประกอบงานปรากฏตั้งแต่ต้น Sprint 3 ไม่ใช่ตอนจบ

\noindent\begin{tabularx}{\dimexpr\textwidth-\leftskip\relax}{@{} L{1.5cm} Y L{2.1cm} @{}}
\toprule
\hrow \thead{กำหนด} & \thead{งาน} & \thead{ผู้รับผิดชอบ} \\
\midrule
28 ส.ค. & รวมงานฝั่ง Back-End ที่เสร็จแล้วเข้าสายหลัก ทั้งรายการอุปกรณ์ การค้นหา ระบบเครดิต วงจรยืม–คืนฝั่งเจ้าหน้าที่ การตรวจสภาพ และการอัปโหลดรูปภาพ พร้อมตัดสินให้เหลือวิธีจัดระดับอุปกรณ์เพียงวิธีเดียว & BE-lead \\
28 ส.ค. & รวมหน้ารับของและหน้าใช้ห้องเข้าสายหลัก & FE-lead \\
29 ส.ค. & Procedure ฝั่งผู้ยืมสำหรับสร้างคำขอยืม ดูคำขอของตนเอง และยกเลิกคำขอ เป็นชิ้นสุดท้ายที่ทำให้วงจรยืม–คืนปิดครบ & BE-dev \\
29 ส.ค. & รวมคีย์ไม่ซ้ำบนอีเมลและรหัสผู้ใช้เข้าสคีมา พร้อมสร้าง Migration ใหม่ขณะที่ฐานข้อมูลยังไม่มีข้อมูลจริง & SA \\
30 ส.ค. & รับไฟล์สัญญาไปใช้แทนสัญญาที่ฝั่ง Front-End เขียนขึ้นเอง และลบไฟล์เดิมทิ้ง & FE-lead \\
31 ส.ค. & เปลี่ยนงานตรวจความสอดคล้องของเวอร์ชันให้ขวางการรวมโค้ด และเพิ่มงานเรียกชุดทดสอบเข้าระบบตรวจสอบอัตโนมัติ & QA-lead \\
1 ก.ย. & ติดตั้งตัวรันชุดทดสอบฝั่ง Front-End & QA-lead \\
2 ก.ย. & ต่อหน้ารายการอุปกรณ์และหน้ารายละเอียดอุปกรณ์เข้ากับข้อมูลจริง & FE-lead \\
3 ก.ย. & ทดสอบวงจรยืม–คืนครบเส้นทางปกติบนสายหลัก ซึ่งเป็นเกณฑ์ที่ยกมาจาก Sprint 2 & PM \\
\bottomrule
\end{tabularx}

\subsection{Sprint 3 ถึง Sprint 5}

\noindent\begin{tabularx}{\dimexpr\textwidth-\leftskip\relax}{@{} L{2.9cm} Y @{}}
\toprule
\hrow \thead{Sprint} & \thead{เป้าหมายและเกณฑ์ ``เสร็จ''} \\
\midrule
\textbf{Sprint 3} \newline 28 ส.ค. – 10 ก.ย. & งานที่ยกมาจาก Sprint 2 · กฎธุรกิจรายระดับ · และการจองสถานที่ (T3) \newline งานที่ยกมาจาก Sprint 2 — รวมงานที่เสร็จแล้วเข้าสายหลัก · สร้างคำขอยืมฝั่งผู้ยืม · ข้อมูลตั้งต้นของอุปกรณ์และห้อง · รับไฟล์สัญญา · เลิกใช้ข้อมูลจำลองในเส้นทางหลัก · ตัวรันชุดทดสอบ (กำหนดรายวันอยู่ในหัวข้อ \ref{sec:next}) \newline กฎรายระดับ — อนุมัติอัตโนมัติสำหรับ T0/T1 · ผู้อนุมัติพิจารณา T2 · จัดเส้นทางอนุมัติตามหน่วยงานเจ้าของ · เลือกและผูกหมายเลขชิ้น (T1–T2) · ตรวจสภาพพร้อมรูปภาพ · ประเมินความเสียหายและระบบเครดิต · แจ้งเตือน \newline การจองสถานที่ T3 — คำนวณและแสดงสล็อตเวลาว่างของห้อง · จองพร้อมกันไม่เกินจำนวนที่กำหนดและกันการจองซ้อนช่วงเวลา · อนุมัติอัตโนมัติตามกฎ IsFree · รอบตรวจสภาพห้องประจำวันของเจ้าหน้าที่ · หน้าจอเลือกช่วงเวลาต่อยอดจากหน้ารายการห้องที่มีอยู่แล้ว \newline เกณฑ์เสร็จ — ยืมอุปกรณ์ข้ามภาควิชาได้ครบวงจรในเส้นทางปกติ · ไม่เหลือข้อมูลจำลองในเส้นทางหลัก · กฎรายระดับและการจองห้องทำงานได้ \newline หมายเหตุ — ตารางห้องและตารางการจองมีอยู่ในฐานข้อมูลแล้ว งาน T3 จึงเป็นการเพิ่ม Procedure และหน้าจอ ไม่ต้องแก้โครงสร้างข้อมูล \\
\textbf{Sprint 4} \newline 11 – 24 ก.ย. & ข้ามหน่วยงานและงานขัดเงา — ตะกร้าพร้อมการแตกอนุมัติตามเจ้าของ · การจองล่วงหน้าถึงสิ้นภาคการศึกษาสำหรับ T1–T3 · การต่ออายุ · การบังคับสิทธิ์ทุกจุดพร้อมบันทึกการตรวจสอบ · หน้าจอที่ยังว่าง · การค้นหา การแบ่งหน้า และการรองรับมือถือ \newline หมายเหตุ — เป็น Sprint สุดท้ายที่เพิ่มความสามารถใหม่ได้ \\
\textbf{Sprint 5} \newline 25 ก.ย. – 1 ต.ค. & หยุดเพิ่มความสามารถและเตรียมนำเสนอ — ทดสอบถดถอยเต็มรูปแบบ · ทดสอบสมรรถนะ · นำระบบขึ้นใช้งานจริง · SRS และ SDS ฉบับทางการ · คู่มือผู้ใช้และคู่มือติดตั้ง · ซ้อมนำเสนอ \\
\bottomrule
\end{tabularx}

\subsection{การประเมินกำหนดเสร็จ}

จุดวัดสามข้อที่ตั้งไว้สำหรับ Sprint 2 ผ่านหนึ่งข้อ คือหน้าจออ่านข้อมูลจากฐานข้อมูลจริงได้แล้ว อีกสองข้อยังไม่ผ่าน คือการรับไฟล์สัญญาไปใช้ และวงจรยืม–คืนที่ทำงานได้จริง

คณะทำงานประเมินว่าโครงการยังเสร็จได้ภายใน 1 ตุลาคม 2569 ตามกำหนดเดิม โดยมีขอบเขตครบตามข้อเสนอโครงการ เพราะส่วนที่ยากที่สุดของ Sprint 2 เขียนเสร็จแล้ว สิ่งที่เหลือคือการรวมงานเข้าด้วยกันซึ่งใช้เวลาน้อยกว่าการเขียนขึ้นใหม่ และ Sprint 3 มีกำลังคนเต็มตลอดช่วง การประเมินนี้ตั้งอยู่บนจุดวัดสามข้อถัดไป ถ้าข้อใดไม่ผ่าน คณะทำงานจะทบทวนขอบเขตของ Sprint 4 ในการทบทวนกลาง Sprint 3

\begin{enumerate}
  \item งานที่เสร็จแล้วทั้งหมดถูกรวมเข้าสายหลักภายใน 29 สิงหาคม 2569
  \item วงจรยืม–คืนเส้นทางปกติทำงานได้จริงบนสายหลักภายใน 3 กันยายน 2569
  \item เส้นทางหลักเลิกใช้ข้อมูลจำลอง และกฎรายระดับกับการจองห้องทำงานได้ ภายในวันสิ้นสุด Sprint 3 คือ 10 กันยายน 2569
\end{enumerate}

% =============================================================================

% =============================================================================
\section*{ภาคผนวก — ค่าที่ใช้พล็อตกราฟ}

ค่าเป้าหมายรายมิติ ณ วันสิ้นสุดของแต่ละ Sprint คูณด้วยน้ำหนักของมิตินั้นแล้วรวมกัน คอลัมน์ ``รวม'' คือค่าที่พล็อตบนเส้นแผนในรูปที่ \ref{fig:scurve}

วันรายงานรอบนี้คือ 27 ส.ค. ตรงกับขอบ Sprint 2 พอดี ค่าแผน \planpct\ จึงอ่านจากแถว ``S2 · 27 ส.ค.'' ในตารางนี้ได้โดยตรง ไม่ต้องสอดแทรกค่า

แถวแรก ``เริ่ม 24 ก.ค.'' คือจุดตั้งต้นของการนับ ทุกมิติเป็นศูนย์ยกเว้น Documentation ที่ 10\% ซึ่งมาจากข้อเสนอโครงการ งานออกแบบที่เสร็จก่อนวันนั้นถือเป็นข้อมูลนำเข้าของงานออกแบบใน Sprint 0 ไม่ใช่ค่าตั้งต้น

\noindent\begin{tabularx}{\dimexpr\textwidth-\leftskip\relax}{@{} L{2.4cm} L{1.0cm} L{1.0cm} L{1.0cm} L{1.0cm} L{1.0cm} L{1.0cm} L{1.0cm} Y @{}}
\toprule
\hrow \thead{ขอบ Sprint} & \thead{} & \thead{BE Des} & \thead{FE Des} & \thead{BE Dev} & \thead{FE Dev} & \thead{Test} & \thead{Doc} & \thead{รวม} \\
\hrow \thead{} & \thead{} & \thead{10\%} & \thead{10\%} & \thead{30\%} & \thead{25\%} & \thead{15\%} & \thead{10\%} & \thead{100\%} \\
\midrule
เริ่ม 24 ก.ค. &  & 0 & 0 & 0 & 0 & 0 & 10 & 1 \\
S0 · 30 ก.ค. &  & 45 & 30 & 5 & 5 & 0 & 20 & 12 \\
S1 · 13 ส.ค. &  & 88 & 75 & 25 & 42 & 10 & 65 & 42 \\
S2 · 27 ส.ค. &  & 95 & 85 & 60 & 65 & 35 & 70 & 65 \\
S3 · 10 ก.ย. &  & 100 & 92 & 78 & 78 & 52 & 76 & 78 \\
S4 · 24 ก.ย. &  & 100 & 100 & 97 & 97 & 80 & 88 & 94 \\
S5 · 1 ต.ค. &  & 100 & 100 & 100 & 100 & 100 & 100 & 100 \\
\bottomrule
\end{tabularx}


\end{document}
